\documentclass[11pt]{article}
\usepackage[utf8]{inputenc}
\usepackage[T1]{fontenc}
\usepackage[margin=0.82in,headheight=15pt]{geometry}
\usepackage{lmodern}
\usepackage{amsmath,amssymb}
\usepackage{array}
\usepackage{booktabs}
\usepackage{float}
\usepackage{graphicx}
\usepackage{longtable}
\usepackage{enumitem}
\usepackage{xcolor}
\usepackage[hidelinks]{hyperref}
\usepackage{titlesec}
\usepackage{fancyhdr}
\graphicspath{{assets/doc_figures/}{doc_figures/}}
\setlength{\parindent}{0pt}
\setlength{\parskip}{0.42em}
\setlist[itemize]{nosep,leftmargin=1.35em}
\setlist[enumerate]{nosep,leftmargin=1.55em}
\definecolor{Accent}{HTML}{0B4F6C}
\titleformat{\section}{\large\bfseries\color{Accent}}{}{0pt}{}
\titleformat{\subsection}{\normalsize\bfseries\color{Accent}}{}{0pt}{}
\titleformat{\subsubsection}{\normalsize\bfseries}{}{0pt}{}
\pagestyle{fancy}
\fancyhf{}
\fancyhead[L]{Chaos Theory Toolbox}
\fancyhead[R]{Manual matematico y tecnico}
\fancyfoot[C]{\thepage}

\newcommand{\code}[1]{\texttt{#1}}
\newcommand{\R}{\mathbb{R}}
\newcommand{\sysfig}[4]{
\begin{figure}[H]
\centering
\includegraphics[width=0.98\linewidth]{systems/#1_phase_timeseries.png}
\caption{#2: (a) atractor o trayectoria visible en el espacio tridimensional usado por la herramienta y (b) series temporales de las variables \(x\), \(y\) y \(z\). Parametros reportados: #3. #4.}
\end{figure}
\begin{figure}[H]
\centering
\includegraphics[width=0.98\linewidth]{systems/#1_projections.png}
\caption{Retratos 2-D de #2: (a) \(x\) contra \(y\), (b) \(x\) contra \(z\) y (c) \(y\) contra \(z\). Parametros reportados: #3. #4.}
\end{figure}
}
\newcommand{\examplefig}[3]{
\begin{figure}[H]
\centering
\includegraphics[width=#2\linewidth]{#1}
\caption{#3}
\end{figure}
}

\begin{document}

{\LARGE\bfseries Diccionario, manual y referencia tecnica}\par
{\small Documento de referencia para la pestana \emph{Diccionario}. El contenido describe el estado actual del codigo, las funciones implementadas, sus limites y la interpretacion matematica de las graficas.}\par
\vspace{0.8em}
\tableofcontents
\newpage

\section{Alcance y limites}

La herramienta es un banco de pruebas para sistemas dinamicos no lineales. Usa PySide6 para la interfaz, Matplotlib/pyqtgraph para graficas y un backend nativo en C para calculos pesados. Permite simular trayectorias, retratos 2D, vista 3D, series temporales, comparacion de integradores, densidad espectral de potencia de Welch, espectro de amplitud, bifurcaciones, autovalores locales, cuencas de clasificacion y exponentes de Lyapunov finitos.

El resultado es numerico y exploratorio. La clase de cuenca ``Acotado residual / no clasificado'' significa que el clasificador rapido no detecto escape, convergencia a equilibrio ni periodicidad simple dentro del horizonte elegido. No afirma que la trayectoria sea caotica.

\begin{longtable}{p{0.18\linewidth}p{0.19\linewidth}p{0.51\linewidth}}
\toprule
Clase & Dimension interna & Alcance actual \\
\midrule
\endhead
Flujos ODE 3D & $(x,y,z)$ & Trayectorias, series, retratos, 3D, bifurcacion, cuencas, equilibrios, autovalores y Lyapunov QR-Benettin. \\
Mapas & 1D o 2D, guardados en 3 columnas & Trayectorias e iteraciones, series visibles y bifurcacion por ultimos valores. Sin cuencas ni Lyapunov en la UI actual. \\
DDE Mackey--Glass & estado escalar con retardo & Se muestran $x(t)$, $x(t-\tau)$ y $\dot{x}(t)$. Sin cuencas ni Lyapunov en la UI actual. \\
Lorenz--96 & $J$ variables, $4\le J\le256$ & Se integra el anillo completo pero se grafican $X_1,X_2,X_3$. \\
\bottomrule
\end{longtable}

\section{Manual de uso}

Windows:
\begin{verbatim}
.\run.ps1
.\run.ps1 -Workers 4
\end{verbatim}
Linux/macOS:
\begin{verbatim}
CHAOS_WORKERS=4 ./run-linux.sh
CHAOS_WORKERS=4 ./run-macos.command
\end{verbatim}

Flujo recomendado:
\begin{enumerate}
  \item seleccionar sistema e integrador;
  \item ajustar parametros, condiciones iniciales, paso \code{dt} y tiempo total \code{T};
  \item simular trayectoria y revisar series;
  \item calcular bifurcacion con una malla pequena antes de aumentar resolucion;
  \item calcular cuenca con baja resolucion, ajustar ventana y despues subir \code{Nx,Ny};
  \item para Lyapunov, aumentar primero el tiempo de quemado y luego el tiempo final;
  \item guardar graficas desde la barra de cada pestana.
\end{enumerate}

\section{Sistemas implementados, ecuaciones, parametros y referencias}

Las ecuaciones se escriben en forma de sistema. En las figuras se usaron los parametros por defecto registrados en \code{SYSTEM\_REGISTRY} y el integrador RK4, salvo mapas discretos donde la iteracion no depende de RK4. Todas las graficas de sistemas usan el mismo color magenta y la misma distribucion: primero trayectoria 3D con series temporales y despues tres retratos planos. En mapas 1D/2D y en DDE, las coordenadas auxiliares que no existen fisicamente pueden ser cero o variables reconstruidas para mantener la misma interfaz visual.

\subsection{Lorenz}
Referencia base: \cite{lorenz1963}.
\[
\left\{\begin{aligned}
\dot{x}&=\sigma(y-x),\\
\dot{y}&=x(\rho-z)-y,\\
\dot{z}&=xy-\beta z.
\end{aligned}\right.
\]
\sysfig{lorenz}{Sistema de Lorenz}{\(\sigma=10,\rho=28,\beta=8/3\)}{Condicion inicial \((0.1,0.1,0.1)\), ancho de paso \(h=0.01\), \(T=40\), \(n=4001\).}

\newpage
\subsection{Rossler}
Referencia base: \cite{rossler1976}.
\[
\left\{\begin{aligned}
\dot{x}&=-y-z,\\
\dot{y}&=x+ay,\\
\dot{z}&=b+z(x-c).
\end{aligned}\right.
\]
\sysfig{rossler}{Sistema de Rossler}{\(a=0.2,b=0.2,c=5.7\)}{Condicion inicial \((0.1,0,0)\), ancho de paso \(h=0.02\), \(T=160\), \(n=8001\).}

\newpage
\subsection{Chua / doble scroll}
Referencias base: \cite{chua1984,chua1986}.
\[
\left\{\begin{aligned}
\dot{x}&=\alpha(y-x-f(x)),\\
\dot{y}&=x-y+z,\\
\dot{z}&=-\beta y,\\
f(x)&=m_1x+\frac{m_0-m_1}{2}\left(|x+1|-|x-1|\right).
\end{aligned}\right.
\]
\sysfig{chua}{Circuito de Chua adimensional}{\(\alpha=15.6,\beta=28,m_0=-1.143,m_1=-0.714\)}{Condicion inicial \((0.1,0,0)\), ancho de paso \(h=0.01\), \(T=80\), \(n=8001\).}

\newpage
\subsection{Chen}
Referencia base: \cite{chen1999}.
\[
\left\{\begin{aligned}
\dot{x}&=a(y-x),\\
\dot{y}&=(c-a)x-xz+cy,\\
\dot{z}&=xy-bz.
\end{aligned}\right.
\]
\sysfig{chen}{Sistema de Chen--Ueta}{\(a=35,b=3,c=28\)}{Condicion inicial \((0.1,0.1,0.1)\), ancho de paso \(h=0.005\), \(T=40\), \(n=8001\).}

\newpage
\subsection{Lu}
Referencia base: \cite{lu2002}.
\[
\left\{\begin{aligned}
\dot{x}&=a(y-x),\\
\dot{y}&=-xz+cy,\\
\dot{z}&=xy-bz.
\end{aligned}\right.
\]
\sysfig{lu}{Sistema de Lu}{\(a=36,b=3,c=20\)}{Condicion inicial \((0.1,0.1,0.1)\), ancho de paso \(h=0.005\), \(T=40\), \(n=8001\).}

\newpage
\subsection{Henon}
Referencia base: \cite{henon1976}.
\[
\left\{\begin{aligned}
x_{n+1}&=1-a x_n^2+y_n,\\
y_{n+1}&=b x_n.
\end{aligned}\right.
\]
\sysfig{henon}{Mapa de Henon}{\(a=1.4,b=0.3\)}{Condicion inicial \((0.1,0.1,0)\), iteracion directa con \(h=1\), \(T=1800\), \(n=1801\).}

\newpage
\subsection{Logistico}
Referencia base: \cite{may1976}.
\[
\left\{\begin{aligned}
x_{n+1}&=r x_n(1-x_n),\\
y_{n+1}&=0,\qquad z_{n+1}=0.
\end{aligned}\right.
\]
\sysfig{logistic}{Mapa logistico}{\(r=3.9\)}{Condicion inicial \((0.2,0,0)\), iteracion directa con \(h=1\), \(T=1200\), \(n=1201\).}

\newpage
\subsection{Ikeda}
Referencia base: \cite{ikeda1979}.
\[
\left\{\begin{aligned}
\theta_n&=0.4-\frac{6}{1+x_n^2+y_n^2},\\
x_{n+1}&=1+u(x_n\cos\theta_n-y_n\sin\theta_n),\\
y_{n+1}&=u(x_n\sin\theta_n+y_n\cos\theta_n).
\end{aligned}\right.
\]
\sysfig{ikeda}{Mapa de Ikeda}{\(u=0.918\)}{Condicion inicial \((0.1,0.1,0)\), iteracion directa con \(h=1\), \(T=1200\), \(n=1201\).}

\newpage
\subsection{Mackey--Glass}
Referencia base: \cite{mackey1977}.
\[
\left\{\begin{aligned}
\dot{x}(t)&=\frac{\beta x(t-\tau)}{1+|x(t-\tau)|^n}-\gamma x(t),\\
y(t)&=x(t-\tau),\\
z(t)&=\dot{x}(t).
\end{aligned}\right.
\]
\sysfig{mackey_glass}{Mackey--Glass}{\(\beta=0.2,\gamma=0.1,n=10,\tau=17\)}{Historia inicial \(x(t)=1.2\), ancho de paso \(h=0.1\), \(T=600\), \(N=6001\).}

\newpage
\subsection{Duffing--Ueda}
Referencias base: \cite{ueda1979,holmes1979}.
\[
\left\{\begin{aligned}
\dot{x}&=y,\\
\dot{y}&=-\delta y-\alpha x-\beta x^3+\gamma\cos\theta,\\
\dot{\theta}&=\omega.
\end{aligned}\right.
\]
\sysfig{duffing_ueda}{Oscilador Duffing--Ueda autonomizado}{\(\delta=0.2,\alpha=-1,\beta=1,\gamma=0.3,\omega=1.2\)}{Condicion inicial \((0.1,0,0)\), ancho de paso \(h=0.01\), \(T=120\), \(n=12001\).}

\newpage
\subsection{Rabinovich--Fabrikant}
Referencia base: \cite{rabinovich1979}.
\[
\left\{\begin{aligned}
\dot{x}&=y(z-1+x^2)+\gamma x,\\
\dot{y}&=x(3z+1-x^2)+\gamma y,\\
\dot{z}&=-2z(\alpha+xy).
\end{aligned}\right.
\]
\sysfig{rabinovich_fabrikant}{Rabinovich--Fabrikant}{\(\alpha=1.1,\gamma=0.87\)}{Condicion inicial \((-1,0,0.5)\), ancho de paso \(h=0.005\), \(T=40\), \(n=8001\).}

\newpage
\subsection{Rikitake}
Referencia base: \cite{rikitake1958}.
\[
\left\{\begin{aligned}
\dot{x}&=-\mu x+yz,\\
\dot{y}&=-\mu y+x(z-a),\\
\dot{z}&=1-xy.
\end{aligned}\right.
\]
\sysfig{rikitake}{Dinamo de Rikitake}{\(\mu=2,a=5\)}{Condicion inicial \((0.1,0.1,0.1)\), ancho de paso \(h=0.01\), \(T=100\), \(n=10001\).}

\newpage
\subsection{Sprott A}
Referencia base: \cite{sprott1994}.
\[
\left\{\begin{aligned}
\dot{x}&=y,\\
\dot{y}&=-x+yz,\\
\dot{z}&=1-y^2.
\end{aligned}\right.
\]
\sysfig{sprott_a}{Flujo Sprott A}{sin parametros}{Condicion inicial \((0.1,0.1,0.1)\), ancho de paso \(h=0.01\), \(T=80\), \(n=8001\).}

\newpage
\subsection{Unified Lorenz--Chen}
Referencia base: \cite{lu2002}.
\[
\left\{\begin{aligned}
\dot{x}&=(25\alpha+10)(y-x),\\
\dot{y}&=(28-35\alpha)x+(29\alpha-1)y-xz,\\
\dot{z}&=-\frac{\alpha+8}{3}z+xy.
\end{aligned}\right.
\]
\sysfig{unified_lorenz_chen}{Sistema unificado de Lorenz-Chen}{\(\alpha=0\)}{Condicion inicial \((0.1,0.1,0.1)\), ancho de paso \(h=0.01\), \(T=40\), \(n=4001\).}

\newpage
\subsection{Sprott B}
Referencia base: \cite{sprott1994}.
\[
\left\{\begin{aligned}
\dot{x}&=yz,\\
\dot{y}&=x-y,\\
\dot{z}&=1-xy.
\end{aligned}\right.
\]
\sysfig{sprott_b}{Flujo Sprott B}{sin parametros}{Condicion inicial \((0.1,0.1,0.1)\), ancho de paso \(h=0.01\), \(T=80\), \(n=8001\).}

\newpage
\subsection{Sprott C}
Referencia base: \cite{sprott1994}.
\[
\left\{\begin{aligned}
\dot{x}&=yz,\\
\dot{y}&=x-y,\\
\dot{z}&=1-x^2.
\end{aligned}\right.
\]
\sysfig{sprott_c}{Flujo Sprott C}{sin parametros}{Condicion inicial \((0.1,0.1,0.1)\), ancho de paso \(h=0.01\), \(T=80\), \(n=8001\).}

\newpage
\subsection{Sprott D}
Referencia base: \cite{sprott1994}.
\[
\left\{\begin{aligned}
\dot{x}&=-y,\\
\dot{y}&=x+z,\\
\dot{z}&=xz+3y^2.
\end{aligned}\right.
\]
\sysfig{sprott_d}{Flujo Sprott D}{sin parametros}{Condicion inicial \((0.1,0.1,0.1)\), ancho de paso \(h=0.01\), \(T=80\), \(n=8001\).}

\newpage
\subsection{Sprott E}
Referencia base: \cite{sprott1994}.
\[
\left\{\begin{aligned}
\dot{x}&=yz,\\
\dot{y}&=x^2-y,\\
\dot{z}&=1-4x.
\end{aligned}\right.
\]
\sysfig{sprott_e}{Flujo Sprott E}{sin parametros}{Condicion inicial \((0.1,0.1,0.1)\), ancho de paso \(h=0.01\), \(T=80\), \(n=8001\).}

\newpage
\subsection{Sprott F}
Referencia base: \cite{sprott1994}.
\[
\left\{\begin{aligned}
\dot{x}&=y+z,\\
\dot{y}&=-x+0.5y,\\
\dot{z}&=x^2-z.
\end{aligned}\right.
\]
\sysfig{sprott_f}{Flujo Sprott F}{sin parametros}{Condicion inicial \((0.1,0.1,0.1)\), ancho de paso \(h=0.01\), \(T=80\), \(n=8001\).}

\newpage
\subsection{Sprott G}
Referencia base: \cite{sprott1994}.
\[
\left\{\begin{aligned}
\dot{x}&=0.4x+z,\\
\dot{y}&=xz-y,\\
\dot{z}&=-x+y.
\end{aligned}\right.
\]
\sysfig{sprott_g}{Flujo Sprott G}{sin parametros}{Condicion inicial \((0.1,0.1,0.1)\), ancho de paso \(h=0.01\), \(T=80\), \(n=8001\).}

\newpage
\subsection{Sprott H}
Referencia base: \cite{sprott1994}.
\[
\left\{\begin{aligned}
\dot{x}&=-y+z^2,\\
\dot{y}&=x+0.5y,\\
\dot{z}&=x-z.
\end{aligned}\right.
\]
\sysfig{sprott_h}{Flujo Sprott H}{sin parametros}{Condicion inicial \((0.1,0.1,0.1)\), ancho de paso \(h=0.01\), \(T=80\), \(n=8001\).}

\newpage
\subsection{Sprott I}
Referencia base: \cite{sprott1994}.
\[
\left\{\begin{aligned}
\dot{x}&=0.2y,\\
\dot{y}&=x+z,\\
\dot{z}&=x+y^2-z.
\end{aligned}\right.
\]
\sysfig{sprott_i}{Flujo Sprott I}{sin parametros}{Condicion inicial \((0.1,0.1,0.1)\), ancho de paso \(h=0.01\), \(T=80\), \(n=8001\).}

\newpage
\subsection{Sprott J}
Referencia base: \cite{sprott1994}.
\[
\left\{\begin{aligned}
\dot{x}&=2z,\\
\dot{y}&=-2y+z,\\
\dot{z}&=-x+y+y^2.
\end{aligned}\right.
\]
\sysfig{sprott_j}{Flujo Sprott J}{sin parametros}{Condicion inicial \((0.1,0.1,0.1)\), ancho de paso \(h=0.01\), \(T=80\), \(n=8001\).}

\newpage
\subsection{Sprott K}
Referencia base: \cite{sprott1994}.
\[
\left\{\begin{aligned}
\dot{x}&=xy-z,\\
\dot{y}&=x-y,\\
\dot{z}&=x+0.3z.
\end{aligned}\right.
\]
\sysfig{sprott_k}{Flujo Sprott K}{sin parametros}{Condicion inicial \((0.1,0.1,0.1)\), ancho de paso \(h=0.01\), \(T=80\), \(n=8001\).}

\newpage
\subsection{Sprott L}
Referencia base: \cite{sprott1994}.
\[
\left\{\begin{aligned}
\dot{x}&=y+3.9z,\\
\dot{y}&=0.9x^2-y,\\
\dot{z}&=1-x.
\end{aligned}\right.
\]
\sysfig{sprott_l}{Flujo Sprott L}{sin parametros}{Condicion inicial \((0.1,0.1,0.1)\), ancho de paso \(h=0.01\), \(T=80\), \(n=8001\).}

\newpage
\subsection{Sprott M}
Referencia base: \cite{sprott1994}.
\[
\left\{\begin{aligned}
\dot{x}&=-z,\\
\dot{y}&=-x^2-y,\\
\dot{z}&=1.7+1.7x+y.
\end{aligned}\right.
\]
\sysfig{sprott_m}{Flujo Sprott M}{sin parametros}{Condicion inicial \((0.1,0.1,0.1)\), ancho de paso \(h=0.01\), \(T=80\), \(n=8001\).}

\newpage
\subsection{Sprott N}
Referencia base: \cite{sprott1994}.
\[
\left\{\begin{aligned}
\dot{x}&=-2y,\\
\dot{y}&=x+z^2,\\
\dot{z}&=1+y-2z.
\end{aligned}\right.
\]
\sysfig{sprott_n}{Flujo Sprott N}{sin parametros}{Condicion inicial \((0.1,0.1,0.1)\), ancho de paso \(h=0.01\), \(T=80\), \(n=8001\).}

\newpage
\subsection{Sprott O}
Referencia base: \cite{sprott1994}.
\[
\left\{\begin{aligned}
\dot{x}&=y,\\
\dot{y}&=x-z,\\
\dot{z}&=x+xz+2.7y.
\end{aligned}\right.
\]
\sysfig{sprott_o}{Flujo Sprott O}{sin parametros}{Condicion inicial \((0.1,0.1,0.1)\), ancho de paso \(h=0.01\), \(T=80\), \(n=8001\).}

\newpage
\subsection{Sprott P}
Referencia base: \cite{sprott1994}.
\[
\left\{\begin{aligned}
\dot{x}&=2.7y+z,\\
\dot{y}&=-x+y^2,\\
\dot{z}&=x+y.
\end{aligned}\right.
\]
\sysfig{sprott_p}{Flujo Sprott P}{sin parametros}{Condicion inicial \((0.1,0.1,0.1)\), ancho de paso \(h=0.01\), \(T=80\), \(n=8001\).}

\newpage
\subsection{Sprott Q}
Referencia base: \cite{sprott1994}.
\[
\left\{\begin{aligned}
\dot{x}&=-z,\\
\dot{y}&=x-y,\\
\dot{z}&=3.1x+y^2+0.5z.
\end{aligned}\right.
\]
\sysfig{sprott_q}{Flujo Sprott Q}{sin parametros}{Condicion inicial \((0.1,0.1,0.1)\), ancho de paso \(h=0.01\), \(T=80\), \(n=8001\).}

\newpage
\subsection{Sprott R}
Referencia base: \cite{sprott1994}.
\[
\left\{\begin{aligned}
\dot{x}&=0.9-y,\\
\dot{y}&=0.4+z,\\
\dot{z}&=xy-z.
\end{aligned}\right.
\]
\sysfig{sprott_r}{Flujo Sprott R}{sin parametros}{Condicion inicial \((0.1,0.1,0.1)\), ancho de paso \(h=0.01\), \(T=80\), \(n=8001\).}

\newpage
\subsection{Sprott S}
Referencia base: \cite{sprott1994}.
\[
\left\{\begin{aligned}
\dot{x}&=x-4y,\\
\dot{y}&=x+z^2,\\
\dot{z}&=1+x.
\end{aligned}\right.
\]
\sysfig{sprott_s}{Flujo Sprott S}{sin parametros}{Condicion inicial \((0.1,0.1,0.1)\), ancho de paso \(h=0.01\), \(T=80\), \(n=8001\).}

\newpage
\subsection{Thomas}
Referencia base: \cite{thomas1999}.
\[
\left\{\begin{aligned}
\dot{x}&=\sin y-bx,\\
\dot{y}&=\sin z-by,\\
\dot{z}&=\sin x-bz.
\end{aligned}\right.
\]
\sysfig{thomas}{Flujo de Thomas}{\(b=0.18\)}{Condicion inicial \((0.1,0,0)\), ancho de paso \(h=0.02\), \(T=140\), \(n=7001\).}

\newpage
\subsection{Hindmarsh--Rose}
Referencia base: \cite{hindmarsh1984}.
\[
\left\{\begin{aligned}
\dot{x}&=y-ax^3+bx^2-z+I,\\
\dot{y}&=c-dx^2-y,\\
\dot{z}&=r(s(x-x_r)-z),\qquad x_r=-1.6.
\end{aligned}\right.
\]
\sysfig{hindmarsh_rose}{Hindmarsh--Rose}{\(a=1,b=3,c=1,d=5,r=0.006,s=4,I=3.25\)}{Condicion inicial \((0.1,0,0)\), ancho de paso \(h=0.02\), \(T=500\), \(n=25001\).}

\newpage
\subsection{Lorenz--96}
Referencia base: \cite{lorenz1996}.
\[
\left\{\begin{aligned}
\dot{X}_j&=(X_{j+1}-X_{j-2})X_{j-1}-X_j+F,\\
j&=0,\ldots,J-1,\qquad X_{j+J}=X_j.
\end{aligned}\right.
\]
\sysfig{lorenz96}{Lorenz--96, primeras tres variables visibles}{\(F=8,J=8\)}{Condicion inicial visible \((X_1,X_2,X_3)=(8.01,8,8)\), ancho de paso \(h=0.01\), \(T=40\), \(n=4001\).}

\newpage
\section{Metodos numericos implementados}

\subsection{Euler explicito}
Euler aproxima el campo vectorial como constante durante un paso:
\[
\left\{\begin{aligned}
k_1&=f(x_n,t_n),\\
x_{n+1}&=x_n+h k_1.
\end{aligned}\right.
\]
Es de orden local \(O(h^2)\) y orden global \(O(h)\). Sirve para inspeccion rapida y para explicar la dinamica, pero en sistemas caoticos puede crear o destruir estructuras por error numerico si \(h\) es grande. Si una trayectoria cambia cualitativamente al pasar de Euler a RK4, la simulacion no debe interpretarse como robusta.

\subsection{Heun / Euler mejorado}
Heun usa un predictor y luego corrige con la pendiente al final del paso:
\[
\left\{\begin{aligned}
k_1&=f(x_n,t_n),\\
\tilde{x}_{n+1}&=x_n+h k_1,\\
k_2&=f(\tilde{x}_{n+1},t_n+h),\\
x_{n+1}&=x_n+\frac{h}{2}(k_1+k_2).
\end{aligned}\right.
\]
Es un Runge--Kutta de orden 2. Reduce sesgo respecto a Euler y permite comparar si la geometria observada depende demasiado del integrador.

\subsection{Runge--Kutta 4}
RK4 evalua cuatro pendientes:
\[
\left\{\begin{aligned}
k_1&=f(x_n,t_n),\\
k_2&=f(x_n+\tfrac{h}{2}k_1,t_n+\tfrac{h}{2}),\\
k_3&=f(x_n+\tfrac{h}{2}k_2,t_n+\tfrac{h}{2}),\\
k_4&=f(x_n+h k_3,t_n+h),\\
x_{n+1}&=x_n+\frac{h}{6}(k_1+2k_2+2k_3+k_4).
\end{aligned}\right.
\]
Es el metodo recomendado por defecto para la exploracion visual actual. No adapta el paso; si hay rigidez, transitorios rapidos o escalas multiples, debe reducirse \(h\) y contrastar con series temporales y diagnosticos.

\section{Funciones y logica metodologica}

\subsection{Simulacion y seleccion de trayectorias}
Para flujos ODE 3D se integra desde \(x(0)=x_0\) hasta \(T\) con \(N=\lfloor T/h\rfloor+1\). La trayectoria numerica es
\[
\mathcal{O}_h(x_0)=\{x_0,x_1,\ldots,x_N\}.
\]
Para mapas se aplica directamente \(x_{n+1}=F(x_n)\). Para Mackey--Glass se guarda una historia constante y se aproxima el retardo por indices discretos. Para Lorenz--96 se integra todo el anillo, pero la UI solo expone tres variables.

Indicaciones: usar primero \(T\) corto para descartar divergencias; despues aumentar \(T\) para ver regimen asintotico. Si se sospecha caos, comparar al menos dos pasos \(h\) y dos integradores.

\examplefig{lorenz_attractor.png}{0.78}{Ejemplo original conservado: atractor de Lorenz con \(\sigma=10,\rho=28,\beta=8/3\). Esta grafica sirve como referencia visual de trayectoria acotada no periodica: la orbita no converge a un punto, sino que visita dos lobulos con alternancia irregular.}

\subsection{Diagramas de bifurcacion}
Un diagrama de bifurcacion muestra como cambia el comportamiento asintotico al variar un parametro \(\mu\). La herramienta barre una malla uniforme
\[
\mu_j=\mu_{\min}+j\frac{\mu_{\max}-\mu_{\min}}{N_\mu-1},\qquad j=0,\ldots,N_\mu-1.
\]
Para cada \(\mu_j\), integra, descarta un transitorio \(T_{\mathrm{trans}}\) y registra eventos durante \(T_{\mathrm{keep}}\).

En Lorenz se usa una seccion de Poincare:
\[
\Sigma=\{(x,y,z):x=0\},\qquad x_{\mathrm{prev}}>0,\;x_{\mathrm{new}}\le 0.
\]
El cruce se interpola linealmente:
\[
\alpha=\frac{x_{\mathrm{prev}}}{x_{\mathrm{prev}}-x_{\mathrm{new}}},\qquad
z_\Sigma=z_{\mathrm{prev}}+\alpha(z_{\mathrm{new}}-z_{\mathrm{prev}}).
\]
Interpretacion: un equilibrio estable no genera nubes de cruces; un ciclo limite produce uno o pocos puntos repetidos en la seccion; una orbita caotica produce una nube extendida. En flujos genericos se usan maximos locales de \(z\):
\[
z_{k-1}>z_{k-2},\qquad z_{k-1}\ge z_k.
\]
En mapas se grafican ultimos valores de \(x_n\).

Una bifurcacion de Hopf ocurre cuando un par complejo conjugado de autovalores cruza el eje imaginario:
\[
\lambda_{1,2}(\mu)=\alpha(\mu)\pm i\omega(\mu),\qquad \alpha(\mu_*)=0,\quad \omega(\mu_*)\ne0.
\]
Alrededor de \(\mu_*\) puede nacer o desaparecer un ciclo limite. En un diagrama por Poincare, un ciclo limite aparece como una rama finita de puntos; duplicaciones de periodo aparecen como desdoblamientos sucesivos; una region caotica aparece como nube de muchos valores.

Indicaciones: empezar con \(N_\mu\approx 200\), \(T_{\mathrm{trans}}\) largo y \(T_{\mathrm{keep}}\) moderado. Si hay saltos o ramas rotas, aumentar transitorio, disminuir \(h\) y probar con/sin continuacion.

\examplefig{logistic_bifurcation.png}{0.84}{Ejemplo original conservado: diagrama de bifurcacion del mapa logistico. Para cada valor de \(r\) se descartan iteraciones transitorias y se grafican valores asintoticos de \(x_n\). Una sola rama indica punto fijo, dos ramas indican periodo 2, duplicaciones sucesivas indican cascada de periodo y las nubes densas corresponden a regimenes caoticos o ventanas periodicas inmersas.}

\subsection{Cuencas y clasificacion de comportamiento}
La cuenca se calcula en el plano
\[
(x(0),y(0),z(0))=(x_0,y_0,z_0^{\mathrm{fijo}}).
\]
Para cada punto se integra y se asigna una clase:
\begin{itemize}
  \item \textbf{Escape}: valores no finitos o radio de escape excedido.
  \item \textbf{Converge a \(E_i\)}: el estado final queda dentro de un radio local del equilibrio \(E_i\).
  \item \textbf{Periodico}: la cola tiene amplitud muy pequena o la seccion de Poincare repite uno o dos grupos estrechos.
  \item \textbf{Acotado residual / no clasificado}: no escapa, no converge y no cumple el criterio periodico rapido dentro del horizonte finito.
\end{itemize}
La convergencia se decide con
\[
k^*=\arg\min_k \|x(T)-E_k\|^2,\qquad \|x(T)-E_{k^*}\|<r_{\mathrm{conv}}.
\]
La periodicidad rapida usa la cola \([T/2,T]\) y agrupa los valores de cruce \(z_\Sigma\). Si no se reconoce periodicidad simple, la trayectoria permanece en la clase residual/no clasificada; no se promueve automaticamente a caos. Esta clasificacion es operacional y rapida; no sustituye un analisis riguroso con exponentes, secciones largas o pruebas de recurrencia.

Indicaciones: comenzar con \(60\times60\), ajustar el plano \(z_0\), despues subir a \(200\times200\) o mas. Para fronteras finas, aumentar \(T_{\mathrm{total}}\) antes de concluir.

\examplefig{lorenz_basin.png}{0.72}{Ejemplo original conservado: cuenca rapida de Lorenz en el plano \(z_0=1\), ventana lejana \([-60,60]\times[-60,60]\). La leyenda representa las clases operacionales disponibles para este caso: acotado residual/no clasificado, converge a \(E_+\) y converge a \(E_-\).}

\examplefig{lorenz_basin_trajectories.png}{0.98}{Ejemplo original conservado: trayectorias iniciadas en regiones distintas de la cuenca de Lorenz. Las condiciones iniciales en zonas rojas y verdes se acercan a vecindades de equilibrios simetricos, mientras que la condicion en la zona azul permanece acotada con recorrido irregular y queda como residual/no clasificada por el clasificador rapido.}

\subsection{Analisis espectral: PSD de Welch y espectro de amplitud}
Despues de retirar el transitorio, el modo recomendado calcula una densidad espectral de potencia (PSD) de Welch unilateral. Cada segmento de longitud \(L\) se centra, se multiplica por una ventana de Hann \(w_k\) y produce
\[
S_m^{(r)}=\frac{h}{\sum_{k=0}^{L-1}w_k^2}
\left|\sum_{k=0}^{L-1}w_k(s_{k+r}-\bar{s}_r)e^{-2\pi i mk/L}\right|^2.
\]
Los bins interiores de frecuencia positiva se duplican y la PSD final promedia los segmentos superpuestos al 50\%. Para una variable con unidades \(U\) y tiempo en segundos, la ordenada tiene unidades \(U^2/\mathrm{Hz}\); no se divide por el maximo espectral.

El modo alternativo conserva las unidades de la variable y muestra la magnitud unilateral del coeficiente de amplitud, normalizada por la suma de la ventana:
\[
A_m=\frac{\left|\sum_{k=0}^{N-1}w_k(s_k-\bar{s})e^{-2\pi i mk/N}\right|}
{\max\!\left(\sum_{k=0}^{N-1}w_k,10^{-300}\right)}.
\]
Ambos modos usan el eje unilateral \(f_m=m/(Nh)\ge0\), sin \texttt{fftshift} ni frecuencias negativas. La PSD se dibuja como curva y el espectro de amplitud con lineas verticales. Si no se fijan limites manuales, la UI recorta el eje alrededor del 98\% central de la medida espectral dominante entre componentes. Los colores siguen la configuracion de \(x(t)\), \(y(t)\) y \(z(t)\).

Un ciclo limite suele producir picos dominantes en frecuencias fundamentales y armonicos. Una trayectoria irregular puede mostrar banda ancha, aunque la interpretacion depende de \(T\), \(h\), ventana temporal y variable observada. Ninguno de los dos espectros certifica por si solo periodicidad, caos, atraccion u ocultamiento.

\examplefig{lorenz_fft.png}{0.78}{Ejemplo de espectro de amplitud para Lorenz con lineas verticales y eje \(x\) unilateral recortado alrededor de las frecuencias dominantes. La presencia de amplitud distribuida en varias frecuencias es consistente con dinamica irregular, pero debe contrastarse con secciones, cuencas y exponentes de Lyapunov.}

\subsection{Exponentes de Lyapunov y metodo de Benettin}
Los exponentes de Lyapunov miden tasas medias de expansion de perturbaciones:
\[
\lambda_i=\lim_{t\to\infty}\frac{1}{t}\log\frac{\|\delta_i(t)\|}{\|\delta_i(0)\|}.
\]
Un exponente maximo positivo indica sensibilidad exponencial promedio a condiciones iniciales. El algoritmo QR-Benettin integra el estado y una base tangente:
\[
\dot{x}=f(x),\qquad \dot{Y}=J(x)Y.
\]
En la implementacion actual:
\[
Y_{n+1}\approx Y_n+hJ(x_n)Y_n,
\]
donde \(J\) se aproxima por diferencias centradas. Cada cierto numero de pasos se aplica
\[
Y=QR,
\]
y se acumulan las diagonales:
\[
\lambda_i(t)\approx \frac{1}{t}\sum_m \log |R_{ii}^{(m)}|.
\]
La reortonormalizacion evita que todas las perturbaciones colapsen hacia la direccion mas expansiva. El resultado depende de \(h\), tiempo de quemado, tiempo final, tolerancia del Jacobiano y frecuencia QR.

Indicaciones: usar \(t_{\mathrm{burn}}\) para retirar transitorios; aumentar \(t_{\mathrm{final}}\) hasta que las curvas \(\lambda_i(t)\) se estabilicen visualmente; repetir con otro \(h\).

\examplefig{lorenz_lyapunov.png}{0.78}{Ejemplo original conservado: estimacion finita de exponentes de Lyapunov para Lorenz por QR-Benettin. La lectura esperada en el regimen clasico es una direccion expansiva, una casi neutra asociada al flujo y una contractiva; la convergencia debe evaluarse por estabilidad temporal, no por un unico valor instantaneo.}

\subsection{Autovalores y clasificacion local}
Para un equilibrio \(x^*\), la linealizacion es
\[
\dot{\xi}=J(x^*)\xi.
\]
Si \(\lambda_i\) son los autovalores:
\[
\left\{\begin{array}{ll}
\Re(\lambda_i)<0\ \forall i & \text{estable localmente},\\
\exists i:\Re(\lambda_i)>0 & \text{inestable},\\
\text{signos mezclados} & \text{silla},\\
\lambda=\alpha\pm i\omega & \text{foco/espiral si }\omega\ne0,\\
\Re(\lambda)=0 & \text{caso no hiperbolico}.
\end{array}\right.
\]
Cuando un par complejo cruza el eje imaginario puede aparecer una bifurcacion de Hopf. La grafica de autovalores no prueba comportamiento global: solo describe la dinamica infinitesimal cerca del equilibrio.

\examplefig{lorenz_spectrum.png}{0.72}{Ejemplo original conservado: espectro de autovalores de los equilibrios de Lorenz. Los puntos a la derecha del eje imaginario indican direcciones locales inestables; pares complejos con parte real distinta de cero indican focos espirales estables o inestables segun el signo de la parte real.}

\begin{thebibliography}{99}
\bibitem{lorenz1963} E. N. Lorenz, ``Deterministic Nonperiodic Flow'', \emph{Journal of the Atmospheric Sciences}, 1963.
\bibitem{rossler1976} O. E. Rossler, ``An equation for continuous chaos'', \emph{Physics Letters A}, 1976.
\bibitem{chua1984} T. Matsumoto, ``A chaotic attractor from Chua's circuit'', \emph{IEEE Transactions on Circuits and Systems}, 1984.
\bibitem{chua1986} L. O. Chua, M. Komuro and T. Matsumoto, ``The double scroll family'', \emph{IEEE Transactions on Circuits and Systems}, 1986.
\bibitem{chen1999} G. Chen and T. Ueta, ``Yet another chaotic attractor'', \emph{International Journal of Bifurcation and Chaos}, 1999.
\bibitem{lu2002} J. Lu and G. Chen, ``A new chaotic attractor coined'', \emph{International Journal of Bifurcation and Chaos}, 2002.
\bibitem{henon1976} M. Henon, ``A two-dimensional mapping with a strange attractor'', \emph{Communications in Mathematical Physics}, 1976.
\bibitem{may1976} R. M. May, ``Simple mathematical models with very complicated dynamics'', \emph{Nature}, 1976.
\bibitem{ikeda1979} K. Ikeda, ``Multiple-valued stationary state and its instability of the transmitted light by a ring cavity system'', \emph{Optics Communications}, 1979.
\bibitem{mackey1977} M. C. Mackey and L. Glass, ``Oscillation and chaos in physiological control systems'', \emph{Science}, 1977.
\bibitem{ueda1979} Y. Ueda, work on randomly transitional phenomena in Duffing-type systems, 1979.
\bibitem{holmes1979} P. Holmes, ``A nonlinear oscillator with a strange attractor'', \emph{Philosophical Transactions of the Royal Society A}, 1979.
\bibitem{rabinovich1979} M. I. Rabinovich and A. L. Fabrikant, ``Stochastic self-modulation of waves in nonequilibrium media'', 1979.
\bibitem{rikitake1958} T. Rikitake, ``Oscillations of a system of disk dynamos'', \emph{Proceedings of the Cambridge Philosophical Society}, 1958.
\bibitem{sprott1994} J. C. Sprott, ``Some simple chaotic flows'', \emph{Physical Review E}, 1994.
\bibitem{thomas1999} R. Thomas, work on deterministic chaos and feedback circuits, 1999.
\bibitem{hindmarsh1984} J. L. Hindmarsh and R. M. Rose, ``A model of neuronal bursting using three coupled first order differential equations'', \emph{Proceedings of the Royal Society B}, 1984.
\bibitem{lorenz1996} E. N. Lorenz, ``Predictability: a problem partly solved'', ECMWF Seminar, 1996.
\bibitem{benettin1980} G. Benettin et al., ``Lyapunov characteristic exponents for smooth dynamical systems'', \emph{Meccanica}, 1980.
\bibitem{guckenheimer1983} J. Guckenheimer and P. Holmes, \emph{Nonlinear Oscillations, Dynamical Systems, and Bifurcations of Vector Fields}, 1983.
\bibitem{kuznetsov1998} Y. A. Kuznetsov, \emph{Elements of Applied Bifurcation Theory}, 1998.
\end{thebibliography}

\vfill
{\small Este PDF se carga en la interfaz cuando \code{QtPdf} esta disponible.}

\end{document}
